\section{Validation, tests et preuves}

\subsection{Stratégie de preuve}

La validation suit le risque : les unités vérifient chaque invariant; les
tests de page reproduisent le DOM, les mutations et la popup; les tests live
chargent réellement l’extension dans Chrome et parcourent une interface
distante. Une vérification manuelle complète les aspects difficiles à
automatiser, notamment zoom, contraste et lecture d’écran.

\begin{figure}[H]
\centering
\begin{tikzpicture}[
  level/.style={trapezium,trapezium left angle=75,trapezium right angle=105,
    draw=Line,minimum height=11mm,align=center,font=\small},
]
\node[level,fill=Mint,minimum width=56mm] (l1) at (0,0)
  {\textbf{Tests live}\\Safinum + UBA dans Chrome};
\node[level,fill=Paper,minimum width=89mm] (l2) at (0,-12mm)
  {\textbf{Tests navigateur}\\popup · mutations · iframe · Shadow DOM};
\node[level,fill=Warm,minimum width=124mm] (l3) at (0,-24mm)
  {\textbf{Tests unitaires et corpus}\\moteur · intégrité · orthographe · modèles · export};
\end{tikzpicture}
\caption{Pyramide de test KORA : précision en bas, parcours réel en haut.}
\end{figure}

\subsection{Résultat automatisé de référence}

Le 14 novembre 2026, la suite cœur a reconstruit les index puis exécuté
27 tests :

\begin{keybox}[Résultat observé]
\texttt{tests 27 · pass 27 · fail 0 · skipped 0}\\
Durée du runner : environ 9,53 s; durée avec précompilation : 15,60 s.\\
Index généré : 504 entrées, 391 clés FR, 390 clés EN.
\end{keybox}

Cette suite couvre notamment les sept erreurs publiques, la parité FR/EN, les
valeurs canoniques, les phrases inédites compositionnelles, l’abstention, les
corrections morphologiques, l’export du modèle et une faute générée pour chaque
libellé éligible.

\subsection{Commandes reproductibles}

\begin{lstlisting}[language=bash]
npm install
npm test                 # moteur, corpus, modèles, export, orthographe
npm run test:spelling    # 1 252 libellés fautifs FR/EN
npm run test:popup       # protocole, thème, activation, restauration
npm run test:pages       # DOM dynamique, attributs, iframe, Shadow DOM
npm run test:live        # parcours Safinum réel dans Chrome
npm run test:uba         # page UBA réelle et Ctrl+A
npm run test:all         # chaîne complète
npm run package          # rebuild, modèle exporté, ZIP < 10 Mio
\end{lstlisting}

\subsection{Matrice des scénarios critiques}

\begin{longtable}{@{}L{43mm}L{55mm}L{57mm}@{}}
\toprule
\textbf{Scénario} & \textbf{Oracle} & \textbf{Couverture}\\
\midrule
\endfirsthead
\toprule
\textbf{Scénario} & \textbf{Oracle} & \textbf{Couverture}\\
\midrule
\endhead
Source française / anglaise &
Même intention source, même sortie \code{fub}. &
Tests moteur et bascule live.\\
Solde et historique &
Montants, signes, dates et devise identiques. &
Tests moteur + tableau de bord live.\\
Calcul des frais &
Montant, frais et total identiques avant/après traduction. &
Tests valeurs + étape frais live.\\
Données bénéficiaire &
Nom, téléphone, motif et note inchangés. &
Tests formulaire, récapitulatif et reçu.\\
Libellé d’un \code{select} &
Texte visible traduit; valeur technique inchangée. &
Test navigateur dédié.\\
Message ajouté après chargement &
Détecté par l’observateur et traité sans hook. &
Ajout, \code{innerHTML}, fragment, remplacement, attribut.\\
Fautes d’orthographe &
Même sortie uniquement si un meilleur sens est unique. &
1\,252 libellés : 628 FR et 624 EN.\\
Texte libre inconnu &
Source conservée; aucune phrase inventée. &
Test de faible confiance.\\
Négation financière inconnue &
Source conservée. &
\texttt{Le transfert n'a pas échoué.}\\
Code technique &
Code conservé; explication visible traduisible. &
HTTP, événements, API et erreurs Safinum.\\
Sélection clavier &
\code{Ctrl+A} sélectionne la page et un champ normalement. &
Tests popup, pages et UBA.\\
Restauration &
Tous les nœuds reviennent à leur texte original. &
Test protocole popup/content script.\\
\bottomrule
\end{longtable}

\subsection{Test live Safinum}

Le scénario automatisé ouvre Chrome avec \code{kora-unpacked}, puis :

\begin{enumerate}
  \item active KORA sur l’environnement Safinum;
  \item vérifie le tableau de bord et les valeurs canoniques;
  \item ajoute après chargement des messages connus, fautifs et variables;
  \item consulte les frais et choisit le type d’opération;
  \item saisit des données fictives de bénéficiaire;
  \item vérifie le récapitulatif et confirme;
  \item contrôle le reçu, la référence et le retour au tableau de bord;
  \item ouvre une seconde page, bascule Safinum en anglais et confirme la
        même sortie \code{fub}.
\end{enumerate}

Le résultat de référence documenté est :

\begin{lstlisting}[language={}]
LIVE PASS: French journey, dynamic content, form values and receipt integrity.
LIVE PASS: English source switched and translated to the same fub interface.
\end{lstlisting}

Le contrôle UBA est facultatif et dépend d’un site tiers. Une exécution
antérieure a vérifié la sélection de page et de champ; lors de la vérification
finale du 14 novembre, la navigation vers ce site a expiré après 90 secondes.
Ce résultat n’est donc pas présenté comme une réussite reconfirmée.

\subsection{Tests du correcteur orthographique}

Le correcteur normalise accents et ponctuation puis autorise de une à quatre
opérations selon la longueur : insertion, suppression, substitution ou
inversion adjacente. Les tests génèrent une variante fautive pour chaque clé
textuelle éligible et exigent exactement la traduction de la clé d’origine.

\begin{tabularx}{\linewidth}{@{}L{39mm}L{50mm}Y@{}}
\toprule
\textbf{Entrée} & \textbf{Sortie attendue} & \textbf{Règle}\\
\midrule
\texttt{sodle insufisant} &
\texttt{Kaalisi yonaa} &
Deux fautes, rapprochement unique.\\
\texttt{insuficient balnce} &
\texttt{Kaalisi yonaa} &
Source anglaise fautive, intention connue.\\
Deux clés également proches &
Texte source &
Ambiguïté = abstention.\\
\texttt{NETWORK\_001} &
\texttt{NETWORK\_001} &
Code protégé, jamais corrigé approximativement.\\
\bottomrule
\end{tabularx}

\subsection{Évaluation du modèle de langue}

La validation neuronale est séparée par phrase : 107 phrases Common Voice ne
participent pas à l’entraînement. Les 262 formulations KORA et 386 candidats
terminologiques enrichissent l’entraînement mais ne gonflent pas cette
métrique. \code{npm run model:demo} vérifie directement les poids livrés :

\begin{lstlisting}[language={}]
transfeer maa -> yettaama, neldaama
kaalisi       -> gondi, yonaa
Votre paiement de 25 000 FCFA a ete refuse.
-> Yobugol maa 25 000 FCFA salaama.
Le transfert n'a pas echoue.
-> original conserve
\end{lstlisting}

Les caractères ASCII de cet extrait simplifient l’affichage du terminal;
l’exécution réelle et les fichiers utilisent l’orthographe UTF-8 complète.

\subsection{Contrôles manuels avant remise}

\begin{itemize}
  \item types d’opération : compte Safinum, Mobile Money, autre banque;
  \item montants : zéro, 25\,000, supérieur au solde, seuils de frais;
  \item erreurs : montant, bénéficiaire, sécurité, réseau et confirmation;
  \item clavier seul, focus visible, zoom 200\,\%, contraste et lecteur d’écran;
  \item activation, suspension, restauration, rechargement et changement
        FR/EN;
  \item vérification visuelle du reçu puis retour au tableau de bord.
\end{itemize}

\begin{warningbox}[Interprétation des preuves]
Les tests établissent l’intégrité du logiciel sur les cas couverts. Ils ne
transforment pas automatiquement chaque proposition terminologique en forme
native validée. La preuve technique et la preuve linguistique restent
séparées et sont toutes deux nécessaires avant production.
\end{warningbox}
